\section{Relatively ample sheaves}
\label{section-relatively-ample}

\noindent
Let $X$ be a scheme and $\mathcal{L}$ an invertible sheaf on $X$.
Then $\mathcal{L}$ is ample on $X$ if $X$ is quasi-compact and
every point of $X$ is contained in an affine open of the form
$X_s$, where $s \in \Gamma(X, \mathcal{L}^{\otimes n})$ and $n \geq 1$, see
Properties, Definition \ref{properties-definition-ample}.
We turn this into a relative notion as follows.

\begin{definition}
\label{definition-relatively-ample}
\begin{reference}
\cite[II Definition 4.6.1]{EGA}
\end{reference}
Let $f : X \to S$ be a morphism of schemes.
Let $\mathcal{L}$ be an invertible $\mathcal{O}_X$-module.
We say $\mathcal{L}$ is {\it relatively ample}, or {\it $f$-relatively ample},
or {\it ample on $X/S$}, or {\it $f$-ample} if $f : X \to S$
is quasi-compact, and if for every affine open $V \subset S$
the restriction of $\mathcal{L}$ to the open subscheme
$f^{-1}(V)$ of $X$ is ample.
\end{definition}

\noindent
We note that the existence of a relatively ample sheaf on $X$ does not
force the morphism $X \to S$ to be of finite type.

\begin{lemma}
\label{lemma-ample-power-ample}
Let $X \to S$ be a morphism of schemes.
Let $\mathcal{L}$ be an invertible $\mathcal{O}_X$-module.
Let $n \geq 1$. Then $\mathcal{L}$ is $f$-ample if and only if
$\mathcal{L}^{\otimes n}$ is $f$-ample.
\end{lemma}

\begin{proof}
This follows from Properties, Lemma \ref{properties-lemma-ample-power-ample}.
\end{proof}

\begin{lemma}
\label{lemma-relatively-ample-separated}
Let $f : X \to S$ be a morphism of schemes.
If there exists an $f$-ample invertible sheaf, then
$f$ is separated.
\end{lemma}

\begin{proof}
Being separated is local on the base (see
Schemes, Lemma \ref{schemes-lemma-characterize-separated} for example;
it also follows easily from the definition).
Hence we may assume $S$ is affine and $X$
has an ample invertible sheaf. In this case the
result follows from
Properties, Lemma \ref{properties-lemma-ample-separated}.
\end{proof}

\noindent
There are many ways to characterize relatively ample invertible
sheaves, analogous to the equivalent conditions in
Properties, Proposition \ref{properties-proposition-characterize-ample}.
We will add these here as needed.

\begin{lemma}
\label{lemma-characterize-relatively-ample}
\begin{reference}
\cite[II, Proposition 4.6.3]{EGA}
\end{reference}
Let $f : X \to S$ be a quasi-compact morphism of schemes.
Let $\mathcal{L}$ be an invertible sheaf on $X$.
The following are equivalent:
\begin{enumerate}
\item The invertible sheaf $\mathcal{L}$ is $f$-ample.
\item There exists an open covering $S = \bigcup V_i$
such that each $\mathcal{L}|_{f^{-1}(V_i)}$ is ample
relative to $f^{-1}(V_i) \to V_i$.
\item There exists an affine open covering $S = \bigcup V_i$
such that each $\mathcal{L}|_{f^{-1}(V_i)}$ is ample.
\item There exists a quasi-coherent graded $\mathcal{O}_S$-algebra
$\mathcal{A}$ and a map of graded $\mathcal{O}_X$-algebras
$\psi : f^*\mathcal{A} \to \bigoplus_{d \geq 0} \mathcal{L}^{\otimes d}$
such that $U(\psi) = X$ and
$$
r_{\mathcal{L}, \psi} :
X
\longrightarrow
\underline{\text{Proj}}_S(\mathcal{A})
$$
is an open immersion (see Constructions, Lemma
\ref{constructions-lemma-invertible-map-into-relative-proj} for notation).
\item The morphism $f$ is quasi-separated and
part (4) above holds with
$\mathcal{A} = f_*(\bigoplus_{d \geq 0} \mathcal{L}^{\otimes d})$
and $\psi$ the adjunction mapping.
\item Same as (4) but just requiring $r_{\mathcal{L}, \psi}$
to be an immersion.
\end{enumerate}
\end{lemma}

\begin{proof}
It is immediate from the definition that (1) implies (2) and
(2) implies (3). It is clear that (5) implies (4).

\medskip\noindent
Assume (3) holds for the affine open covering $S = \bigcup V_i$.
We are going to show (5) holds.
Since each $f^{-1}(V_i)$ has an ample invertible sheaf we see
that $f^{-1}(V_i)$ is separated
(Properties, Lemma \ref{properties-lemma-ample-separated}).
Hence $f$ is separated. By
Schemes, Lemma \ref{schemes-lemma-push-forward-quasi-coherent}
we see that $\mathcal{A} = f_*(\bigoplus_{d \geq 0} \mathcal{L}^{\otimes d})$
is a quasi-coherent graded $\mathcal{O}_S$-algebra.
Denote $\psi : f^*\mathcal{A} \to \bigoplus_{d \geq 0} \mathcal{L}^{\otimes d}$
the adjunction mapping.
The description of the open $U(\psi)$ in
Constructions, Section
\ref{constructions-section-invertible-relative-proj}
and the definition of ampleness
of $\mathcal{L}|_{f^{-1}(V_i)}$ show that $U(\psi) = X$.
Moreover, Constructions,
Lemma \ref{constructions-lemma-invertible-map-into-relative-proj} part (3)
shows that the restriction of $r_{\mathcal{L}, \psi}$ to
$f^{-1}(V_i)$ is the same as the morphism from
Properties, Lemma \ref{properties-lemma-map-into-proj}
which is an open immersion according to
Properties, Lemma \ref{properties-lemma-ample-immersion-into-proj}.
Hence (5) holds.

\medskip\noindent
Let us show that (4) implies (1). Assume (4).
Denote $\pi : \underline{\text{Proj}}_S(\mathcal{A}) \to S$
the structure morphism. Choose $V \subset S$ affine open. By
Constructions, Definition \ref{constructions-definition-relative-proj}
we see that $\pi^{-1}(V) \subset \underline{\text{Proj}}_S(\mathcal{A})$
is equal to $\text{Proj}(A)$ where $A = \mathcal{A}(V)$
as a graded ring. Hence $r_{\mathcal{L}, \psi}$ maps
$f^{-1}(V)$ isomorphically onto
a quasi-compact open of $\text{Proj}(A)$.
Moreover, $\mathcal{L}^{\otimes d}$ is isomorphic to the pullback of
$\mathcal{O}_{\text{Proj}(A)}(d)$ for some $d \geq 1$.
(See part (3) of Constructions,
Lemma \ref{constructions-lemma-invertible-map-into-relative-proj}
and the final statement of Constructions,
Lemma \ref{constructions-lemma-invertible-map-into-proj}.)
This implies that $\mathcal{L}|_{f^{-1}(V)}$ is ample by
Properties, Lemmas \ref{properties-lemma-open-in-proj-ample}
and \ref{properties-lemma-ample-power-ample}.

\medskip\noindent
Assume (6). By the equivalence of (1) - (5) above we see that the
property of being relatively ample on $X/S$ is local on $S$. Hence
we may assume that $S$ is affine, and we have to show that
$\mathcal{L}$ is ample on $X$. In this case the morphism
$r_{\mathcal{L}, \psi}$ is identified with the morphism, also denoted
$r_{\mathcal{L}, \psi} : X \to \text{Proj}(A)$ associated to the map
$\psi : A = \mathcal{A}(V) \to \Gamma_*(X, \mathcal{L})$.
(See references above.) As above we also see that
$\mathcal{L}^{\otimes d}$ is the pullback of the sheaf
$\mathcal{O}_{\text{Proj}(A)}(d)$ for some $d \geq 1$.
Moreover, since $X$ is quasi-compact we
see that $X$ gets identified with a closed subscheme of a
quasi-compact open subscheme $Y \subset \text{Proj}(A)$.
By
Constructions, Lemma
\ref{constructions-lemma-ample-on-proj}
(see also
Properties, Lemma
\ref{properties-lemma-open-in-proj-ample})
we see that $\mathcal{O}_Y(d')$ is an ample invertible sheaf on
$Y$ for some $d' \geq 1$. Since the restriction of an ample
sheaf to a closed subscheme is ample, see
Properties, Lemma
\ref{properties-lemma-ample-on-closed}
we conclude that the pullback of
$\mathcal{O}_Y(d')$ is ample. Combining these results with
Properties, Lemma
\ref{properties-lemma-ample-power-ample}
we conclude that $\mathcal{L}$ is ample as desired.
\end{proof}

\begin{lemma}
\label{lemma-ample-over-affine}
\begin{reference}
\cite[II Corollary 4.6.6]{EGA}
\end{reference}
Let $f : X \to S$ be a morphism of schemes.
Let $\mathcal{L}$ be an invertible $\mathcal{O}_X$-module.
Assume $S$ affine.
Then $\mathcal{L}$ is $f$-relatively ample if and only
if $\mathcal{L}$ is ample on $X$.
\end{lemma}

\begin{proof}
Immediate from Lemma \ref{lemma-characterize-relatively-ample}
and the definitions.
\end{proof}

\begin{lemma}
\label{lemma-quasi-affine-O-ample}
\begin{reference}
\cite[II Proposition 5.1.6]{EGA}
\end{reference}
Let $f : X \to S$ be a morphism of schemes. Then $f$ is quasi-affine
if and only if $\mathcal{O}_X$ is $f$-relatively ample.
\end{lemma}

\begin{proof}
Follows from Properties, Lemma \ref{properties-lemma-quasi-affine-O-ample}
and the definitions.
\end{proof}

\begin{lemma}
\label{lemma-pullback-ample-tensor-relatively-ample}
Let $f : X \to Y$ be a morphism of schemes, $\mathcal{M}$
an invertible $\mathcal{O}_Y$-module, and $\mathcal{L}$ an
invertible $\mathcal{O}_X$-module.
\begin{enumerate}
\item If $\mathcal{L}$ is $f$-ample and $\mathcal{M}$
is ample, then $\mathcal{L} \otimes f^*\mathcal{M}^{\otimes a}$ is ample
for $a \gg 0$.
\item If $\mathcal{M}$ is ample
and $f$ quasi-affine, then $f^*\mathcal{M}$ is ample.
\end{enumerate}
\end{lemma}

\begin{proof}
Assume $\mathcal{L}$ is $f$-ample and $\mathcal{M}$ ample.
By assumption $Y$ and $f$ are quasi-compact (see
Definition \ref{definition-relatively-ample} and
Properties, Definition \ref{properties-definition-ample}).
Hence $X$ is quasi-compact. By Properties, Lemma
\ref{properties-lemma-ample-separated} the scheme $Y$ is
separated and by Lemma \ref{lemma-relatively-ample-separated}
the morphism $f$ is separated. Hence $X$ is separated by
Schemes, Lemma \ref{schemes-lemma-separated-permanence}.
Pick $x \in X$. We can choose $m \geq 1$
and $t \in \Gamma(Y, \mathcal{M}^{\otimes m})$ such that $Y_t$
is affine and $f(x) \in Y_t$. Since $\mathcal{L}$ restricts to an
ample invertible sheaf on $f^{-1}(Y_t) = X_{f^*t}$
we can choose $n \geq 1$ and $s \in \Gamma(X_{f^*t}, \mathcal{L}^{\otimes n})$
with $x \in (X_{f^*t})_s$ with $(X_{f^*t})_s$ affine.
By Properties, Lemma \ref{properties-lemma-invert-s-sections} part (2)
whose assumptions are satisfied by the above,
there exists an integer $e \geq 1$ and a section
$s' \in \Gamma(X, \mathcal{L}^{\otimes n} \otimes f^*\mathcal{M}^{\otimes em})$
which restricts to $s(f^*t)^e$ on $X_{f^*t}$. For any $b > 0$
consider the section $s'' = s'(f^*t)^b$ of
$\mathcal{L}^{\otimes n} \otimes f^*\mathcal{M}^{\otimes (e + b)m}$.
Then $X_{s''} = (X_{f^*t})_s$ is an affine open of $X$ containing $x$.
Picking $b$ such that $n$ divides $e + b$ we see
$\mathcal{L}^{\otimes n} \otimes f^*\mathcal{M}^{\otimes (e + b)m}$
is the $n$th power of $\mathcal{L} \otimes f^*\mathcal{M}^{\otimes a}$
for some $a$ and we can get any $a$ divisible by $m$ and big enough.
Since $X$ is quasi-compact a finite number of these affine opens
cover $X$. We conclude that for some $a$ sufficiently divisible and
large enough the invertible sheaf
$\mathcal{L} \otimes f^*\mathcal{M}^{\otimes a}$ is ample on $X$.
On the other hand, we know that $\mathcal{M}^{\otimes c}$
(and hence its pullback to $X$) is globally generated for all $c \gg 0$
by Properties, Proposition \ref{properties-proposition-characterize-ample}.
Thus $\mathcal{L} \otimes f^*\mathcal{M}^{\otimes a + c}$ is ample
(Properties, Lemma \ref{properties-lemma-ample-tensor-globally-generated})
for $c \gg 0$ and (1) is proved.

\medskip\noindent
Part (2) follows from Lemma \ref{lemma-quasi-affine-O-ample},
Properties, Lemma \ref{properties-lemma-ample-power-ample}, and
part (1).
\end{proof}

\begin{lemma}
\label{lemma-ample-composition}
Let $g : Y \to S$ and $f : X \to Y$ be morphisms of schemes.
Let $\mathcal{M}$ be an invertible $\mathcal{O}_Y$-module.
Let $\mathcal{L}$ be an invertible $\mathcal{O}_X$-module.
If $S$ is quasi-compact, $\mathcal{M}$ is $g$-ample, and
$\mathcal{L}$ is $f$-ample, then
$\mathcal{L} \otimes f^*\mathcal{M}^{\otimes a}$
is $g \circ f$-ample for $a \gg 0$.
\end{lemma}

\begin{proof}
Let $S = \bigcup_{i = 1, \ldots, n} V_i$ be a finite affine open covering.
By Lemma \ref{lemma-characterize-relatively-ample}
it suffices to prove that 
$\mathcal{L} \otimes f^*\mathcal{M}^{\otimes a}$
is ample on $(g \circ f)^{-1}(V_i)$ for $i = 1, \ldots, n$.
Thus the lemma follows from
Lemma \ref{lemma-pullback-ample-tensor-relatively-ample}.
\end{proof}

\begin{lemma}
\label{lemma-ample-base-change}
Let $f : X \to S$ be a morphism of schemes.
Let $\mathcal{L}$ be an invertible $\mathcal{O}_X$-module.
Let $S' \to S$ be a morphism of schemes.
Let $f' : X' \to S'$ be the base change of $f$ and denote
$\mathcal{L}'$ the pullback of $\mathcal{L}$ to $X'$.
If $\mathcal{L}$ is $f$-ample, then $\mathcal{L}'$ is $f'$-ample.
\end{lemma}

\begin{proof}
By Lemma \ref{lemma-characterize-relatively-ample} it suffices
to find an affine open covering $S' = \bigcup U'_i$
such that $\mathcal{L}'$ restricts to an ample invertible
sheaf on $(f')^{-1}(U_i')$ for all $i$. We may choose $U'_i$
mapping into an affine open $U_i \subset S$. In this case the
morphism $(f')^{-1}(U'_i) \to f^{-1}(U_i)$ is affine as a base
change of the affine morphism $U'_i \to U_i$
(Lemma \ref{lemma-base-change-affine}). Thus
$\mathcal{L}'|_{(f')^{-1}(U'_i)}$ is ample by
Lemma \ref{lemma-pullback-ample-tensor-relatively-ample}.
\end{proof}

\begin{lemma}
\label{lemma-ample-permanence}
Let $g : Y \to S$ and $f : X \to Y$ be morphisms of schemes.
Let $\mathcal{L}$ be an invertible $\mathcal{O}_X$-module.
If $\mathcal{L}$ is $g \circ f$-ample and $f$ is
quasi-compact\footnote{This follows if $g$ is quasi-separated by
Schemes, Lemma \ref{schemes-lemma-quasi-compact-permanence}.}
then $\mathcal{L}$ is $f$-ample.
\end{lemma}

\begin{proof}
Assume $f$ is quasi-compact and $\mathcal{L}$ is $g \circ f$-ample.
Let $U \subset S$ be an affine open and let $V \subset Y$ be
an affine open with $g(V) \subset U$.
Then $\mathcal{L}|_{(g \circ f)^{-1}(U)}$ is ample on
$(g \circ f)^{-1}(U)$ by assumption.
Since $f^{-1}(V) \subset (g \circ f)^{-1}(U)$ we see that
$\mathcal{L}|_{f^{-1}(V)}$ is ample on $f^{-1}(V)$ by
Properties, Lemma \ref{properties-lemma-ample-on-locally-closed}.
Namely, $f^{-1}(V) \to (g \circ f)^{-1}(U)$ is a quasi-compact
open immersion by
Schemes, Lemma \ref{schemes-lemma-quasi-compact-permanence}
as $(g \circ f)^{-1}(U)$ is separated
(Properties, Lemma \ref{properties-lemma-ample-separated})
and $f^{-1}(V)$ is quasi-compact (as $f$ is quasi-compact).
Thus we conclude that $\mathcal{L}$ is $f$-ample by
Lemma \ref{lemma-characterize-relatively-ample}.
\end{proof}







